\section*{Chapitre \ref{ch:b1:isostasy-gravity} --- Isostasie et anomalies de pesanteur}

\begin{solution}{exo:b1:isostasy-gravity:1}
(a) $\qty{-150}{\milli\gal} = -150\times\qty{e-5}{m.s^{-2}} = \qty{-1.5e-3}{m.s^{-2}}$,
et $\num{1.5e-3}/9.81 = \num{1.53e-4}$ de $g$. (b) $\qty{1}{\milli\gal} =
\qty{e-5}{m.s^{-2}}$, \num{1.02e-6} de $g$~; $\qty{1}{\micro\gal} =
\qty{e-8}{m.s^{-2}}$, \num{1.02e-9} de $g$~: un gravimètre qui résout un
microgal mesure la pesanteur à environ une partie sur un milliard.
\end{solution}

\begin{solution}{exo:b1:isostasy-gravity:2}
(a) À l'équateur les deux sinus s'annulent~: $\gamma = \qty{9.780327}{m.s^{-2}}$. À
\ang{45}, $\sin^2\varphi = 0.5$ et $\sin^2 2\varphi = 1$~:
$\gamma = 9.780327\times(1 + 0.0026512 - 0.0000058) = \qty{9.806200}{m.s^{-2}}$.
Aux pôles, $\sin^2\varphi = 1$ et $\sin 2\varphi = 0$~:
$\gamma = 9.780327\times1.0053024 = \qty{9.832186}{m.s^{-2}}$.
(b) $9.832186 - 9.780327 = \qty{0.051859}{m.s^{-2}}$, c'est-à-dire
\qty{5186}{\milli\gal}. (c) $3392/5186 = 0.65$~: l'effet centrifuge explique
environ les deux tiers de la différence~; l'aplatissement, qui rapproche les pôles
du centre, le reste.
\end{solution}

\begin{solution}{exo:b1:isostasy-gravity:3}
$r = h\rho_c/(\rho_m - \rho_c) = 3\times2800/500 = \qty{16.8}{km}$. La \omterm{def:b1:layered-earth:layers}{croûte} est
plus épaisse de $h + r = \qty{19.8}{km}$ (c'est-à-dire $h\rho_m/(\rho_m - \rho_c) =
3\times3300/500$).
\end{solution}

\begin{solution}{exo:b1:isostasy-gravity:4}
(a) $1/0.3086 = \qty{3.24}{m}$. (b) $0.01/0.3086 = \qty{0.032}{m}$~: les hauteurs
doivent être connues à environ \qty{3}{cm}, c'est pourquoi un levé enregistre la
hauteur de chaque station par positionnement par satellite ou par nivellement.
\end{solution}

\begin{solution}{exo:b1:isostasy-gravity:5}
(a) $\rho = 2700\times100/(100 + h)$~: \qty{2673}{kg.m^{-3}} pour \qty{1}{km},
\qty{2621}{kg.m^{-3}} pour \qty{3}{km}, \qty{2571}{kg.m^{-3}} pour \qty{5}{km}.
(b) $\rho = (2700\times100 - 1030\times4)/96 = \num{265880}/96 =
\qty{2770}{kg.m^{-3}}$~: la colonne sous la mer est plus dense. (c) Le modèle de Pratt
demande des changements latéraux de masse volumique à profondeur constante. Sous le
plancher océanique de tels changements existent~: la \omterm{def:b1:layered-earth:lithosphere}{lithosphère} jeune et chaude près
d'une dorsale est plus légère que la \omterm{def:b1:layered-earth:lithosphere}{lithosphère} ancienne et froide loin d'elle
(\cref{ch:b1:heat-budget}). Sous les continents, au contraire, on voit l'épaisseur de
la \omterm{def:b1:layered-earth:layers}{croûte} changer avec le relief.
\end{solution}

\begin{solution}{exo:b1:isostasy-gravity:6}
(a) Station 1~: $g - \gamma = 9.807751 - 9.808911 = \qty{-0.001160}{m.s^{-2}} =
\qty{-116.0}{\milli\gal}$~; \omterm{def:b1:isostasy-gravity:anomalies}{correction à l'air libre} $0.3086\times350 =
\qty{108.0}{\milli\gal}$, donc $\Delta g_{\mathrm{FA}} = \qty{-8.0}{\milli\gal}$~;
\omterm{def:b1:isostasy-gravity:anomalies}{correction de Bouguer} $0.112\times350 = \qty{39.2}{\milli\gal}$, donc
$\Delta g_{\mathrm B} = \qty{-47.2}{\milli\gal}$.
Station 2~: $g - \gamma = \qty{-535.9}{\milli\gal}$~; \omterm{def:b1:isostasy-gravity:anomalies}{correction à l'air libre}
$0.3086\times1850 = \qty{570.9}{\milli\gal}$, donc $\Delta g_{\mathrm{FA}} =
\qty{+35.0}{\milli\gal}$~; \omterm{def:b1:isostasy-gravity:anomalies}{correction de Bouguer} $0.112\times1850 =
\qty{207.2}{\milli\gal}$, donc $\Delta g_{\mathrm B} = \qty{-172.2}{\milli\gal}$.
(b) À la station 2 l'\omterm{def:b1:isostasy-gravity:anomalies}{anomalie de Bouguer} est fortement négative~: sous les montagnes
il y a un déficit de masse, une racine de \omterm{def:b1:layered-earth:layers}{croûte} légère. L'\omterm{def:b1:isostasy-gravity:anomalies}{anomalie à l'air libre}, petite
devant la \omterm{def:b1:isostasy-gravity:anomalies}{correction de Bouguer}, montre que le déficit équilibre presque
les montagnes~: elles sont largement compensées (à environ $172/207$, soit
\qty{83}{\percent}, si le modèle de la plaque vaut pour toute la chaîne).
\end{solution}

\begin{solution}{exo:b1:isostasy-gravity:7}
(a) Comme dans la démonstration de la proposition, une mince couche à la distance $z$
contribue $G\rho\,2\pi z\,h\int_0^a s\,(s^2 + z^2)^{-3/2}\,\mathrm ds$, l'intégrale
étant maintenant arrêtée en $s = a$~:
$\int_0^a s\,(s^2 + z^2)^{-3/2}\,\mathrm ds = 1/z - 1/\sqrt{z^2 + a^2}$. D'où
$g = 2\pi G\rho h\,\bigl(1 - z/\sqrt{z^2 + a^2}\bigr)$, qui tend vers la valeur de la
plaque quand $a \to \infty$. (b) $z/\sqrt{z^2 + a^2} = 0.1$ donne
$\sqrt{z^2 + a^2} = 10z$, $a = \sqrt{99}\,z = \qty{9.95}{km}$ pour
$z = \qty{1}{km}$. (c) La roche comprise dans un rayon d'environ dix fois sa profondeur
sous la station fournit \qty{90}{\percent} de l'attraction~: la correction de terrain
demande le relief sur un rayon de plusieurs kilomètres à dizaines de kilomètres
autour de chaque station, et pas seulement sous elle.
\end{solution}

\begin{solution}{exo:b1:isostasy-gravity:8}
(a) $r' = d(\rho_c - \rho_w)/(\rho_m - \rho_c) = 4\times1770/500 =
\qty{14.16}{km}$, et la \omterm{def:b1:layered-earth:layers}{croûte} fait $35 - 4 - 14.16 = \qty{16.84}{km}$ d'épaisseur.
(b) Le modèle d'Airy à une seule masse volumique de \omterm{def:b1:layered-earth:layers}{croûte} prédit une \omterm{def:b1:layered-earth:layers}{croûte} océanique plus de
trois fois trop épaisse. La \omterm{def:b1:layered-earth:layers}{croûte} océanique n'est pas de la \omterm{def:b1:layered-earth:layers}{croûte} continentale amincie~: c'est
une roche différente, plus dense, et les deux niveaux de la Terre sont deux types de
\omterm{def:b1:layered-earth:layers}{croûte}, comme le suggère la \omterm{def:b1:isostasy-gravity:hypsometry}{courbe hypsométrique}.
\end{solution}

\begin{solution}{exo:b1:isostasy-gravity:9}
(a) Une seconde d'arc vaut $\pi/(180\times3600) = 1/\num{206265}$ de radian~:
\qty{5.236}{''} $= \qty{2.538e-5}{rad}$ et \qty{15.885}{''} $=
\qty{7.701e-5}{rad}$. (b) $a_h = g\theta$~: $9.79\times\num{2.538e-5} =
\qty{2.485e-4}{m.s^{-2}} = \qty{24.9}{\milli\gal}$ observés, et
$9.79\times\num{7.701e-5} = \qty{75.4}{\milli\gal}$ calculés à partir des montagnes
visibles. (c) $1 - 5.236/15.885 = 0.670$~: la compensation a annulé environ les deux
tiers de l'attraction des montagnes visibles.
\end{solution}

\begin{solution}{exo:b1:isostasy-gravity:10}
(a) À l'extérieur, une sphère attire comme une masse ponctuelle en son centre, avec
$GM/(x^2 + z^2)$ le long de la droite vers le centre~; la composante verticale est ceci
multiplié par le cosinus $z/\sqrt{x^2 + z^2}$, ce qui donne $GMz/(x^2 + z^2)^{3/2}$.
(b) $M = \tfrac43\pi(\qty{2000}{m})^3\times(\qty{-200}{kg.m^{-3}}) =
\qty{-6.70e12}{kg}$, et à l'aplomb du centre $GM/z^2 = \num{6.674e-11}\times
(\num{-6.70e12})/(4000)^2 = \qty{-2.80e-5}{m.s^{-2}} = \qty{-2.80}{\milli\gal}$.
(c) $(1 + x^2/z^2)^{-3/2} = 1/2$ donne $x = z\sqrt{2^{2/3} - 1} = 0.766\,z =
\qty{3.07}{km}$. (d) La demi-largeur à mi-hauteur du profil donne la
profondeur~: $z = 1.305\,x_{1/2}$. Un corps plus profond donne une anomalie plus large et plus plate.
\end{solution}

\begin{solution}{exo:b1:isostasy-gravity:11}
(a) $V = \tfrac13\pi R^2 H = \tfrac13\pi(\qty{15e3}{m})^2(\qty{3e3}{m}) =
\qty{7.07e11}{m^3}$, and $\Delta m = (2800 - 1030)V = \qty{1.25e15}{kg}$.
(b) Le centre de masse se trouve à \qty{0.75}{km} au-dessus du fond, soit à
$d = 5 - 0.75 = \qty{4.25}{km}$ sous la surface de la mer. L'anomalie de potentiel
y est $T = -G\Delta m/d$, donc $N = G\Delta m/(\gamma d) =
\num{6.674e-11}\times\num{1.25e15}/(9.81\times4250) = \qty{2.0}{m}$.
(c) L'essentiel de la masse du cône se trouve dans sa large partie inférieure, à 5 à
\qty{16}{km} du point au-dessus du sommet~: la moyenne de $1/r$ sur la masse
est plus petite que $1/\qty{4.25}{km}$, et la masse ponctuelle surestime le
potentiel. En sommant de minces disques on obtient environ \qty{1.1}{m}.
(d) Une racine de \omterm{def:b1:layered-earth:layers}{croûte} légère en dessous mettrait un déficit de même masse plus
profond~; les deux potentiels s'annulent presque, et la bosse serait bien plus petite,
mais pas nulle, puisque le déficit est plus loin de la surface que l'excès.
\end{solution}

\begin{solution}{exo:b1:isostasy-gravity:12}
(a) La vitesse de soulèvement est $-\mathrm dw/\mathrm dt = w_0\,e^{-t/\tau}/\tau = w/\tau$.
(b) $w = 818\,e^{-10000/4000} = 818\times0.0821 = \qty{67.1}{m}$, et la vitesse est
$\qty{67.1}{m}/\num{4000}\,\unit{\year} = \qty{16.8}{mm.\year^{-1}}$, plus d'une fois et
demie la valeur mesurée. (c) Avec $\tau = 3100$ ans, $w =
818\,e^{-3.226} = \qty{32.5}{m}$ et la vitesse est de \qty{10.5}{mm.an^{-1}}~; avec
$\tau = \num{68000}$ ans, $w = 818\,e^{-0.147} = \qty{706}{m}$ et la vitesse
est de \qty{10.4}{mm.an^{-1}}. Les deux impliquent des dépressions restantes très différentes.
Une dépression non compensée est un déficit de masse de roche du \omterm{def:b1:layered-earth:layers}{manteau}~: sur une vaste
région elle donne une \omterm{def:b1:isostasy-gravity:anomalies}{anomalie à l'air libre} d'environ $-2\pi G\rho_m w$, c'est-à-dire
$-0.0419\times3.3\times706 \approx \qty{-98}{\milli\gal}$ pour le cas lent,
contre \qty{-4}{\milli\gal} pour le cas rapide. Un levé gravimétrique (ou la hauteur
des anciens rivages) permet de les distinguer. (d) La constante de temps mesure la vitesse à laquelle
le \omterm{def:b1:layered-earth:layers}{manteau} reflue sous la dépression~: sa viscosité.
\end{solution}

\begin{solution}{pb:b1:isostasy-gravity:1}
\textbf{1.} $r = h\rho_c/(\rho_m - \rho_c) = 4\times2700/600 = \qty{18}{km}$.

\textbf{2.} La \omterm{def:b1:layered-earth:layers}{croûte} fait $35 + 4 + 18 = \qty{57}{km}$ d'épaisseur~; sa base, le \omterm{def:b1:layered-earth:discontinuity}{Moho},
se trouve à $35 + 18 = \qty{53}{km}$ sous le niveau de la mer.

\textbf{3.} Plateau~: \qty{57}{km} de \omterm{def:b1:layered-earth:layers}{croûte} et $60 - 53 = \qty{7}{km}$ de \omterm{def:b1:layered-earth:layers}{manteau},
$\num{57000}\times2700 + \num{7000}\times3300 = \num{1.539e8} + \num{2.31e7} =
\qty{1.770e8}{kg.m^{-2}}$. Plaine~: $\num{35000}\times2700 + \num{25000}\times3300
= \num{9.45e7} + \num{8.25e7} = \qty{1.770e8}{kg.m^{-2}}$. Elles sont égales, comme
l'exige la \cref{prop:b1:isostasy-gravity:equal-mass}.

\textbf{4.} $p = g\times\text{masse par unité d'aire} = 9.81\times\num{1.770e8} =
\qty{1.74e9}{Pa}$, environ \qty{1.7}{GPa}, sous les deux colonnes.

\textbf{5.} La racine doit compenser, par sa plus faible masse volumique, la masse du
plateau~: chaque mètre de racine retire $\rho_m - \rho_c = \qty{600}{kg.m^{-3}}$ de
masse, et chaque mètre de plateau ajoute \qty{2700}{kg.m^{-3}}, donc $2700/600 = 4.5$
mètres de racine sont nécessaires. Avec une \omterm{def:b1:layered-earth:layers}{croûte} de \qty{2900}{kg.m^{-3}}, le rapport
serait $2900/400 = 7.25$~: plus le contraste est faible, plus la racine est profonde.

\textbf{6.} On trouverait le \omterm{def:b1:layered-earth:discontinuity}{Moho}, la forte augmentation de la vitesse sismique,
s'approfondir sous le bord du plateau, de \qty{35}{km} sous la plaine
à environ \qty{53}{km} sous le plateau.

\textbf{7.} $\rho = 2700\times100/104 = \qty{2596}{kg.m^{-3}}$.

\textbf{8.} $\num{104000}\times2596.2 = \qty{2.700e8}{kg.m^{-2}}$ et
$\num{100000}\times2700 = \qty{2.700e8}{kg.m^{-2}}$.

\textbf{9.} $2700 - 2596 = \qty{104}{kg.m^{-3}}$, soit $104/2700 =
\qty{3.85}{\percent}$.

\textbf{10.} $\Delta\rho/\rho = \alpha\,\Delta T$ donne $\Delta T =
0.0385/\num{3e-5} = \qty{1280}{K}$ en moyenne sur toute la colonne de \qty{104}{km}.
Cela ferait fondre une bonne part de la colonne~: la chaleur seule ne peut rendre ce
plateau assez léger, et le mécanisme de Pratt ne peut pas être le principal ici.

\textbf{11.} Les colonnes de Pratt n'ont pas de racine~: le \omterm{def:b1:layered-earth:discontinuity}{Moho} reste à environ \qty{35}{km}
sous le plateau, comme sous la plaine.

\textbf{12.} \qty{52}{km} est proche des \qty{53}{km} d'Airy et loin des \qty{35}{km} de
Pratt~: le plateau est compensé surtout par une racine crustale.

\textbf{13.} $0.0419\times2.7\times4000 = \qty{452.5}{\milli\gal}$.

\textbf{14.} Non compensé (\cref{prop:b1:isostasy-gravity:broad})~:
$\Delta g_{\mathrm{FA}} \approx \qty{+452.5}{\milli\gal}$,
$\Delta g_{\mathrm B} \approx 0$.

\textbf{15.} Compensé~: $\Delta g_{\mathrm{FA}} \approx 0$ et
$\Delta g_{\mathrm B} \approx \qty{-452.5}{\milli\gal}$. La masse manquante de la racine
donne $-0.0419\times0.6\times\num{18000} = \qty{-452.5}{\milli\gal}$, la même chose,
puisque $\Delta\rho\,r = \rho_c h$.

\textbf{16.} $\sin^2\ang{30} = 0.25$ et $\sin^2\ang{60} = 0.75$~:
$\gamma = 9.780327\times(1 + 0.0013256 - 0.0000044) = \qty{9.793249}{m.s^{-2}}$.
Puis $g - \gamma = 9.781085 - 9.793249 = \qty{-1216.4}{\milli\gal}$~; la \omterm{def:b1:isostasy-gravity:anomalies}{correction à
l'air libre} est $0.3086\times4000 = \qty{1234.4}{\milli\gal}$, donc
$\Delta g_{\mathrm{FA}} = \qty{+18.0}{\milli\gal}$, et
$\Delta g_{\mathrm B} = 18.0 - 452.5 = \qty{-434.5}{\milli\gal}$.

\textbf{17.} $434.5/452.5 = 0.96$~: le plateau est compensé à environ \qty{96}{\percent}
(la petite anomalie positive à l'air libre est la part non compensée de
\qty{4}{\percent}).

\textbf{18.} Non. Près du centre d'un large plateau chaque couche anormale agit
comme une plaque, dont l'attraction ne dépend pas de sa profondeur~: seul compte le déficit
de masse par unité d'aire, et il est le même, $\rho_c h$, dans les deux modèles. Les deux
ne diffèrent que près des bords, où la profondeur des masses compensatrices change
la forme de l'anomalie.

\textbf{19.} Loin du bord de part et d'autre l'\omterm{def:b1:isostasy-gravity:anomalies}{anomalie à l'air libre} est proche de zéro.
À quelques kilomètres à l'intérieur du bord, le plateau sous la station est proche et
ne s'étend loin que d'un côté~: il tire encore avec presque toute sa valeur de plaque.
La racine, à \qty{35}{km} ou plus en dessous, est vue de là comme un corps qui se termine
presque à la verticale~: elle tire avec à peine plus de la moitié de son effet. La
différence est un pic positif. À quelques kilomètres à l'extérieur, le plateau tire
très peu, alors que la racine tire encore avec presque la moitié de son effet~: un creux
négatif. Le profil est une paire positif--négatif de part et d'autre de chaque bord, comme sur la
figure du chapitre.

\textbf{20.} La \omterm{def:b1:layered-earth:layers}{croûte} fait maintenant $57 - 1 = \qty{56}{km}$ d'épaisseur, soit $\qty{21}{km}$ de plus
que celle de la plaine. À l'équilibre cet excès se partage en $h' + r'$ avec
$r' = 4.5\,h'$~: $h' = 21/5.5 = \qty{3.818}{km}$ et $r' = \qty{17.18}{km}$.

\textbf{21.} Le plateau se tient maintenant à \qty{3818}{m}.

\textbf{22.} La roche maintenant en surface était \qty{1}{km} sous l'ancienne surface,
à \qty{3000}{m}~; elle est maintenant à \qty{3818}{m}~: elle est remontée de \qty{818}{m},
c'est-à-dire $1000\times\rho_c/\rho_m$.

\textbf{23.} Chaque kilomètre érodé abaisse la surface de $1 - 2700/3300 =
0.182$ de kilomètre, donc la ramener vers le bas de \qty{4}{km} demande
$4/0.182 = \qty{22}{km}$ d'érosion, toute l'épaisseur en excès $h + r$. Dans une
vieille chaîne arasée, les roches en surface étaient donc autrefois à une vingtaine
de kilomètres de profondeur~: les racines profondes de la chaîne sont exhumées.

\textbf{24.} La surface s'abaisse de $1000 - 818$~: pour chaque kilomètre de roche
érodé, l'abaissement net de la surface est de
\textbf{\qty[text-series-to-math=true]{182}{m}}.
\end{solution}
